% ##############################################################################
\section{From probe results to verdicts}
\label{sec:verdict}

% ==============================================================================
\subsection{Evidence classification}
\label{sec:evidence}

The checker reads request rows of a contract window and classifies each row
twice, once for latency and once for quality. A row is \emph{eligible} for a
metric, with a success or a failure, or it is \emph{excluded}, or it reveals a
\emph{checker error}. Table~\ref{tab:evidence} lists the quality rules.
Latency rules are similar, except that canary rows are excluded from latency by
default, because short synthetic prompts would bias the buyer-experienced
latency downward.

\begin{table}[htbp]
  \centering
  \small
  \begin{tabularx}{\linewidth}{@{}p{0.34\linewidth}p{0.27\linewidth}X@{}}
    \toprule
    Row & Classification & Reason code \\
    \midrule
    Canary, status \texttt{ok}, grade 1 & eligible, success &
      \texttt{ELIGIBLE\_CORRECT} \\
    Canary, status \texttt{ok}, grade 0 & eligible, failure &
      \texttt{ELIGIBLE\_INCORRECT} \\
    Canary, timeout or upstream error & eligible, failure &
      \texttt{ELIGIBLE\_SELLER\_FAILURE} \\
    Canary, buyer error, gateway error, or spend-cap refusal & excluded &
      \texttt{EXCLUDED\_BUYER\_ERROR}, \texttt{EXCLUDED\_GATEWAY\_ERROR},
      \texttt{EXCLUDED\_REFUSED\_CAP} \\
    Canary answered by a provider other than the pinned one & excluded &
      \texttt{EXCLUDED\_ATTRIBUTION\_MISMATCH} \\
    Canary, status \texttt{ok}, grade missing & checker error &
      \texttt{MISSING\_CANARY\_GRADE} \\
    Buyer request that carries a grade & checker error &
      \texttt{UNEXPECTED\_CANARY\_GRADE} \\
    Any row with an unknown status & checker error &
      \texttt{UNKNOWN\_REQUEST\_STATUS} \\
    Buyer request & not applicable to quality & \texttt{NOT\_APPLICABLE} \\
    \bottomrule
  \end{tabularx}
  \caption{Quality classification of request rows (\texttt{classify\_quality}).}
  \label{tab:evidence}
\end{table}

The classification implements property P5. Only a provider timeout, an upstream
error, or a wrong answer from the pinned provider counts against the seller.
A single checker error, such as a grade that is missing on a successful
response, makes the whole evidence set unhealthy, and the verdict becomes
\texttt{INSUFFICIENT} with the disposition \texttt{CHECKER\_UNHEALTHY}. The
design therefore prefers withholding a verdict to blaming a seller for a
checker fault.

% ==============================================================================
\subsection{Bound, rule, and verdict}
\label{sec:verdict_rule}

For a metric with $s$ successes among $n$ eligible rows, the checker computes a
one-sided upper confidence bound $U(s, n)$ at confidence $1-\delta$ with
$\delta = 0.05$. For $n \geq 10$ it uses the Wilson bound
\begin{equation}
  \label{eq:wilson}
  U_{\text{W}}(s, n) = \min\!\left\{1,\
  \frac{\hat p + \dfrac{z^2}{2n} + z\sqrt{\dfrac{\hat p(1-\hat p)}{n} +
  \dfrac{z^2}{4n^2}}}{1 + \dfrac{z^2}{n}}\right\},
  \qquad \hat p = \frac{s}{n},\quad z = \Phi^{-1}(1-\delta),
\end{equation}
following Wilson \cite{wilson1927} and the comparison of interval methods by
Brown, Cai, and DasGupta \cite{brown2001}. For $n < 10$ it uses the exact
Clopper--Pearson bound \cite{clopper1934}, which is the value $\theta$ with
$\Pr[\mathrm{Bin}(n, \theta) \leq s] = \delta$, and $U = 1$ when $s = n$. A metric with promised rate $\rho$ (here
$\rminattr$ for latency and $\qminattr$ for quality) and minimum evidence
$n_{\min}$ is a \emph{conclusive failure} when
\begin{equation}
  \label{eq:conclusive_failure}
  n \geq n_{\min} \;\wedge\; U(s, n) < \rho.
\end{equation}
The minima are 20 buyer requests for latency and 8 canaries for quality. The
test is the one-sided hypothesis test of the white paper with the null
hypothesis ``the true rate is at least $\rho$'': the null is rejected only if
even the optimistic end of the interval lies below the promise.

The verdict applies the following precedence, shown in
Figure~\ref{fig:verdict-flow}. An unhealthy evidence set gives
\texttt{INSUFFICIENT}. Otherwise a conclusive failure of either metric gives
\texttt{FAILED}, with a reason code for each failing metric. Otherwise a missing
minimum on either metric gives \texttt{INSUFFICIENT}. Otherwise the verdict is
\texttt{PASSED}. A conclusive failure of one metric therefore takes precedence
over missing evidence on the other, and passing requires enough evidence for
both. Because the two metrics are tested separately at the same
level, the probability that a seller exactly at both promises receives a
failing verdict is at most approximately $2\delta$ by a union bound.

% rendered_images:begin
% ```graphviz
% digraph VerdictDecision {
%   // process : blue | PASSED : sage green | FAILED : coral | INSUFFICIENT : orange | decision : white
%   bgcolor="transparent";
%   pad="0.10";
%   splines=spline;
%   nodesep=0.25;
%   ranksep=0.38;
%   node [shape=box, style="rounded,filled", fillcolor="#FFFFFF", color="#9AA9B8", penwidth=1.8, fontname="Helvetica", fontcolor="#243B53", fontsize=11, margin="0.14,0.08", height=0.42];
%   edge [color="#A3B1C0", penwidth=1.3, arrowhead=vee, arrowsize=0.75, fontname="Helvetica", fontsize=9.5, fontcolor="#7B8794"];
%   rankdir=TB;
%   Ev  [label=<<b>Window evidence</b><br/><font point-size="9" color="#6B4517">buyer rows and canary rows</font>>, fillcolor="#FBEBD4", color="#D9A85F", fontcolor="#6B4517"];
%   Cl  [label=<<b>Classify every row</b><br/><font point-size="9" color="#41719C">eligible, excluded, or checker error</font>>, fillcolor="#D3E3F3", color="#7CA6CE", fontcolor="#1F4E79"];
%   D1  [label="Any checker error?"];
%   Un  [label=<<b>INSUFFICIENT</b><br/><font point-size="9" color="#6B4517">CHECKER_UNHEALTHY</font>>, fillcolor="#FBEBD4", color="#D9A85F", fontcolor="#6B4517"];
%   As  [label=<<b>Assess latency and quality</b><br/><font point-size="9" color="#41719C">successes s, total n, upper bound U</font>>, fillcolor="#D3E3F3", color="#7CA6CE", fontcolor="#1F4E79"];
%   D2  [label=<Some metric has<br/>n ≥ minimum and U &lt; promise?>];
%   Fa  [label=<<b>FAILED</b><br/><font point-size="9" color="#6B2A2A">LATENCY_BELOW_PROMISE and/or<br/>QUALITY_BELOW_PROMISE</font>>, fillcolor="#F6C6C6", color="#D98C8C", fontcolor="#6B2A2A"];
%   D3  [label=<Both metrics have<br/>n ≥ minimum?>];
%   In  [label=<<b>INSUFFICIENT</b><br/><font point-size="9" color="#6B4517">INSUFFICIENT_LATENCY_EVIDENCE and/or<br/>INSUFFICIENT_QUALITY_EVIDENCE</font>>, fillcolor="#FBEBD4", color="#D9A85F", fontcolor="#6B4517"];
%   Pa  [label=<<b>PASSED</b><br/><font point-size="9" color="#4F7A5A">COMPLIANT</font>>, fillcolor="#DFEDE0", color="#8FB79A", fontcolor="#2E5A3D"];
% 
%   Ev -> Cl -> D1;
%   D1 -> Un [label="Yes"];
%   D1 -> As [label="No"];
%   As -> D2;
%   D2 -> Fa [label="Yes"];
%   D2 -> D3 [label="No"];
%   D3 -> In [label="No"];
%   D3 -> Pa [label="Yes"];
%   { rank=same; D1; Un; }
%   { rank=same; D2; Fa; }
%   { rank=same; D3; In; }
% }
% ```
% label=fig:verdict-flow
% caption=Verdict decision flow. Blue is a process, orange marks evidence and insufficient outcomes, coral marks failure, and green marks a pass.
% rendered_images:end
% render_images:begin
\begin{figure}[H]
  \includegraphics[width=\linewidth]{07_verdict.tex.figs/07_verdict.1.png}
  \caption{Verdict decision flow. Blue is a process, orange marks evidence and insufficient outcomes, coral marks failure, and green marks a pass.}
  \label{fig:verdict-flow}
\end{figure}
% render_images:end

% ==============================================================================
\subsection{Comparison with the lower-bound rule of the white paper}
\label{sec:lower_vs_upper}

The white paper flags a contract when the \emph{lower} confidence bound on
the success rate is below $\rminattr$. The prototype's design notes record the
change to an upper bound as decision D-1 and describe it as proposed to the
white paper's author, so it is an open design choice rather than an agreed
revision. The two rules differ in what they treat as evidence of a violation.
The lower bound lies below the point estimate by construction, so a seller whose
true rate equals or exceeds the promise has a high probability of a lower bound
below the promise on a small window. The upper bound lies above the point
estimate, so a flag requires that the data are inconsistent with
compliance at the chosen confidence.

Table~\ref{tab:rules} compares the probability of a flag under both rules by
exact binomial enumeration. Both columns use the one-sided 95\% Wilson interval,
so that the table isolates the direction of the bound; all scenarios have
$n \geq 10$. The scenarios restate those of the design notes, and the values
agree with the script \texttt{scripts/detection\_math.py}.

\begin{table}[htbp]
  \centering
  \small
  \begin{tabularx}{\linewidth}{@{}Xrr@{}}
    \toprule
    Scenario (true rate vs.\ promise, $n$ eligible rows) & Lower-bound rule &
      Upper-bound rule \\
    \midrule
    Healthy canaries, 0.98 vs.\ 0.85, $n = 10$ & 1.000 & 0.00003 \\
    Healthy canaries, 0.95 vs.\ 0.85, $n = 10$ & 1.000 & 0.001 \\
    Weaker model, canaries 0.70 vs.\ 0.85, $n = 10$ & 1.000 & 0.350 \\
    Weaker model, canaries 0.50 vs.\ 0.85, $n = 10$ & 1.000 & 0.828 \\
    Healthy latency, 0.97 vs.\ 0.90, $n = 50$ & 0.445 & 0.00002 \\
    Healthy latency, 0.95 vs.\ 0.90, $n = 50$ & 0.721 & 0.0008 \\
    Slowed seller, latency 0.20 vs.\ 0.90, $n = 50$ & 1.000 & 1.000 \\
    \bottomrule
  \end{tabularx}
  \caption{Probability that a contract is flagged under the lower-bound and the
  upper-bound rule.}
  \label{tab:rules}
\end{table}

Under the lower-bound rule every scenario with ten canaries is flagged with
probability one, including compliant sellers, because even ten successes out of
ten give a Wilson lower bound of about $0.79$, which is below $0.85$. The
upper-bound rule flags compliant sellers rarely and still flags a seller whose
canary accuracy is $0.5$ in about 83\% of contracts.

The upper-bound rule is not exact. When the true rate equals the promise, so
that the null hypothesis holds with equality, the probability of a flag with the
promise $0.85$ is $0.021$ for $n = 8$ (exact bound), $0.050$ for $n = 10$, and
$0.067$ for $n = 20$ and $n = 40$ (Wilson bound). The Wilson bound is therefore
slightly liberal for these sizes, and the nominal 5\% level is approximate.

\subsection{Separating latency and quality}
\label{sec:separating}

A merged success indicator, as in \eqref{eq:wp_success}, dilutes quality
failures when canaries are a small fraction of the requests. Consider a window
of 50 buyer requests that all succeed and 10 canaries on which a substituted
model has accuracy $0.5$. The merged success rate is about $0.917$, against a
promise of $0.90$, and the merged test with 60 requests flags this contract with
probability $0.026$. Testing the ten canaries separately against
$\qminattr = 0.85$ flags it with probability $0.828$. This is the argument for
factoring the indicator in Section~\ref{sec:background_refinements}.

% ==============================================================================
\subsection{Detection power}
\label{sec:detection_power}

Let $a$ be the true canary accuracy of a seller on the active bank. The
probability that the quality metric is a conclusive failure with $n$ usable
canaries is
\begin{equation}
  \label{eq:power}
  \Pr[\text{FAILED}] = \sum_{s=0}^{n} \binom{n}{s} a^{s} (1-a)^{n-s}\,
    \mathbb{1}\!\left[U(s, n) < \qminattr\right].
\end{equation}
Figure~\ref{fig:detection-power} evaluates \eqref{eq:power} with
$\qminattr = 0.85$ for $n$ from 8 to 40, and Table~\ref{tab:power} lists the
smallest $n$ with $\Pr[\text{FAILED}] \geq 0.8$.

\begin{figure}[htbp]
  \centering
  \includegraphics[width=\linewidth]{figures/detection_power.png}
  \caption{Probability of a failing quality verdict as a function of the number
  of usable canaries, for several true accuracies $a$ and a promise of 0.85.}
  \label{fig:detection-power}
\end{figure}

\begin{table}[htbp]
  \centering
  \small
  \begin{tabular}{@{}lrrrrrr@{}}
    \toprule
    True accuracy $a$ & 0.8 & 0.75 & 0.7 & 0.6 & 0.5 & 0.4 \\
    \midrule
    Smallest $n$ with $\Pr[\text{FAILED}] \geq 0.8$ & 334 & 84 & 40 & 16 & 10 & 8 \\
    \bottomrule
  \end{tabular}
  \caption{Evidence needed for 80\% detection probability at a promise of 0.85
  and 95\% one-sided confidence.}
  \label{tab:power}
\end{table}

Four observations follow. First, the curves are not monotone in $n$. The
number of failures that the bound tolerates increases at particular values of
$n$, which produces the sawtooth pattern, and the rule switches from the exact
to the Wilson bound at $n = 10$. Second, power depends on the margin between
$a$ and the promise. With the MVP defaults ($n = 10$, promise $0.85$), a seller
at $a = 0.5$ is flagged with probability $0.83$, and with probability $0.97$ in
at least one of two independent contracts. A seller at $a = 0.7$ is flagged with
probability $0.35$, and a seller at $a = 0.8$ with probability $0.12$.
Third, a seller whose true accuracy lies slightly below the promise is detected
only with large samples: $a = 0.8$ needs $n = 334$ for 80\% power. Fourth, a
seller with $a \geq 0.9$ complies with the promise and is flagged rarely, as the
test intends.

The relevant accuracy is the accuracy of the \emph{substitute} on the
\emph{calibrated} bank. Calibration constrains a single weak model to at most
one correct attempt in three on every retained canary. A different substitute,
for example a quantized copy of the contracted model, can score much higher on
the retained canaries and is correspondingly harder to detect. The analysis in
this section therefore describes the statistical behavior of the rule given $a$
and does not estimate $a$ for any real substitute.
